%%%%%%%%%%%%%%%%%%%%%%%%%%%%%%%%%%%%%%%%
%
% $ Autor: Gruppe 07 $
% $ Projekt: Magic Faucet Fountain $
% $ Pfad: LaTeXTemplate/Nano33BLESense/Contents/General/Materialliste Magic Faucet Fountain.tex $
% $ Kurzbeschreibung: Dieses Dokument zeigt die Materialliste. $
% !TeX spellcheck = de_DE/GB
% !TeX program = pdflatex
% !BIB program = biber/bibtex
% !TeX encoding = utf8
%
%%%%%%%%%%%%%%%%%%%%%%%%%%%%%%%%%%%%%%%%


\chapter{Softwareliste Magic Faucet Fountain}

\begin{longtable}{p{3cm}p{2cm}p{5cm}p{4cm}}
	\textbf{Name} & \textbf{Version} & \textbf{Beschreibung} & \textbf{Link} \\
	\hline
	Arduino IDE 
	& 2.3.6 
	& Die Arduino IDE ermöglicht es, auf dem Computer einen funktionsfähigen Code für seinen Arduino zu schreiben, zu kompilieren und auf den Arduino zu übertragen.
	& \href{https://www.arduino.cc/}{Arduino Homepage} \\
	
	\hline
	nRF Connect
	& 2.7.19
	& Mit der nRF Connect-App kann man via BLE mit einem nRF BLE Modul kommunizieren. Ganz einfach mit einer Smartphone App.
	& \href{https://www.nordicsemi.com/}{nRF Homepage}\\
	
	\hline
	SolidWorks
	& Jahr 2024
	& SolidWorks ist ein sehr vielfältiges CAD-Programm. Es dient zum Beispiel zur erstellung von 3D-Druck teilen und die darauffolgende Weiterverwendung in STL-Datein.
	& \href{https://www.solidworks.com/de}{SolidWorks Homepage}\\
	
	\hline
	Prusa Slicer
	& 2.9.2
	& Mit dem Prusa Slicer kann man STL Datein in GCode für Prusa 3D-Drucker umwandeln.
	& \href{https://www.prusa3d.com/de/page/prusaslicer_424/}{Prusa Slicer}\\
	\hline

\end{longtable}

\textbf{Stand: 19.06.2025}
